\documentclass[a4paper,12pt,twoside]{memoir}

% Castellano
\usepackage[spanish,es-tabla]{babel}
\selectlanguage{spanish}
\usepackage[utf8]{inputenc}
\usepackage[T1]{fontenc}
\usepackage{lmodern} % Scalable font
\usepackage{microtype}
\usepackage{placeins}

\RequirePackage{booktabs}
\RequirePackage[table]{xcolor}
\RequirePackage{xtab}
\RequirePackage{multirow}

% Links
\PassOptionsToPackage{hyphens}{url}\usepackage[colorlinks]{hyperref}
\hypersetup{
	allcolors = {red}
}

% Ecuaciones
\usepackage{amsmath}

% Rutas de fichero / paquete
\newcommand{\ruta}[1]{{\sffamily #1}}

% Párrafos
\nonzeroparskip

% Huérfanas y viudas
\widowpenalty100000
\clubpenalty100000

% Imágenes

% Comando para insertar una imagen en un lugar concreto.
% Los parámetros son:
% 1 --> Ruta absoluta/relativa de la figura
% 2 --> Texto a pie de figura
% 3 --> Tamaño en tanto por uno relativo al ancho de página
\usepackage{graphicx}
\newcommand{\imagen}[3]{
	\begin{figure}[!h]
		\centering
		\includegraphics[width=#3\textwidth]{#1}
		\caption{#2}\label{fig:#1}
	\end{figure}
	\FloatBarrier
}

% Comando para insertar una imagen sin posición.
% Los parámetros son:
% 1 --> Ruta absoluta/relativa de la figura
% 2 --> Texto a pie de figura
% 3 --> Tamaño en tanto por uno relativo al ancho de página
\newcommand{\imagenflotante}[3]{
	\begin{figure}
		\centering
		\includegraphics[width=#3\textwidth]{#1}
		\caption{#2}\label{fig:#1}
	\end{figure}
}

% El comando \figura nos permite insertar figuras comodamente, y utilizando
% siempre el mismo formato. Los parametros son:
% 1 --> Porcentaje del ancho de página que ocupará la figura (de 0 a 1)
% 2 --> Fichero de la imagen
% 3 --> Texto a pie de imagen
% 4 --> Etiqueta (label) para referencias
% 5 --> Opciones que queramos pasarle al \includegraphics
% 6 --> Opciones de posicionamiento a pasarle a \begin{figure}
\newcommand{\figuraConPosicion}[6]{%
  \setlength{\anchoFloat}{#1\textwidth}%
  \addtolength{\anchoFloat}{-4\fboxsep}%
  \setlength{\anchoFigura}{\anchoFloat}%
  \begin{figure}[#6]
    \begin{center}%
      \Ovalbox{%
        \begin{minipage}{\anchoFloat}%
          \begin{center}%
            \includegraphics[width=\anchoFigura,#5]{#2}%
            \caption{#3}%
            \label{#4}%
          \end{center}%
        \end{minipage}
      }%
    \end{center}%
  \end{figure}%
}

%
% Comando para incluir imágenes en formato apaisado (sin marco).
\newcommand{\figuraApaisadaSinMarco}[5]{%
  \begin{figure}%
    \begin{center}%
    \includegraphics[angle=90,height=#1\textheight,#5]{#2}%
    \caption{#3}%
    \label{#4}%
    \end{center}%
  \end{figure}%
}
% Para las tablas
\newcommand{\otoprule}{\midrule [\heavyrulewidth]}
%
% Nuevo comando para tablas pequeñas (menos de una página).
\newcommand{\tablaSmall}[5]{%
 \begin{table}
  \begin{center}
   \rowcolors {2}{gray!35}{}
   \begin{tabular}{#2}
    \toprule
    #4
    \otoprule
    #5
    \bottomrule
   \end{tabular}
   \caption{#1}
   \label{tabla:#3}
  \end{center}
 \end{table}
}

%
% Nuevo comando para tablas pequeñas (menos de una página).
\newcommand{\tablaSmallSinColores}[5]{%
 \begin{table}[H]
  \begin{center}
   \begin{tabular}{#2}
    \toprule
    #4
    \otoprule
    #5
    \bottomrule
   \end{tabular}
   \caption{#1}
   \label{tabla:#3}
  \end{center}
 \end{table}
}

\newcommand{\tablaApaisadaSmall}[5]{%
\begin{landscape}
  \begin{table}
   \begin{center}
    \rowcolors {2}{gray!35}{}
    \begin{tabular}{#2}
     \toprule
     #4
     \otoprule
     #5
     \bottomrule
    \end{tabular}
    \caption{#1}
    \label{tabla:#3}
   \end{center}
  \end{table}
\end{landscape}
}

%
% Nuevo comando para tablas grandes con cabecera y filas alternas coloreadas en gris.
\newcommand{\tabla}[6]{%
  \begin{center}
    \tablefirsthead{
      \toprule
      #5
      \otoprule
    }
    \tablehead{
      \multicolumn{#3}{l}{\small\sl continúa desde la página anterior}\\
      \toprule
      #5
      \otoprule
    }
    \tabletail{
      \hline
      \multicolumn{#3}{r}{\small\sl continúa en la página siguiente}\\
    }
    \tablelasttail{
      \hline
    }
    \bottomcaption{#1}
    \rowcolors {2}{gray!35}{}
    \begin{xtabular}{#2}
      #6
      \bottomrule
    \end{xtabular}
    \label{tabla:#4}
  \end{center}
}

%
% Nuevo comando para tablas grandes con cabecera.
\newcommand{\tablaSinColores}[6]{%
  \begin{center}
    \tablefirsthead{
      \toprule
      #5
      \otoprule
    }
    \tablehead{
      \multicolumn{#3}{l}{\small\sl continúa desde la página anterior}\\
      \toprule
      #5
      \otoprule
    }
    \tabletail{
      \hline
      \multicolumn{#3}{r}{\small\sl continúa en la página siguiente}\\
    }
    \tablelasttail{
      \hline
    }
    \bottomcaption{#1}
    \begin{xtabular}{#2}
      #6
      \bottomrule
    \end{xtabular}
    \label{tabla:#4}
  \end{center}
}

%
% Nuevo comando para tablas grandes sin cabecera.
\newcommand{\tablaSinCabecera}[5]{%
  \begin{center}
    \tablefirsthead{
      \toprule
    }
    \tablehead{
      \multicolumn{#3}{l}{\small\sl continúa desde la página anterior}\\
      \hline
    }
    \tabletail{
      \hline
      \multicolumn{#3}{r}{\small\sl continúa en la página siguiente}\\
    }
    \tablelasttail{
      \hline
    }
    \bottomcaption{#1}
  \begin{xtabular}{#2}
    #5
   \bottomrule
  \end{xtabular}
  \label{tabla:#4}
  \end{center}
}



\definecolor{cgoLight}{HTML}{EEEEEE}
\definecolor{cgoExtralight}{HTML}{FFFFFF}

%
% Nuevo comando para tablas grandes sin cabecera.
\newcommand{\tablaSinCabeceraConBandas}[5]{%
  \begin{center}
    \tablefirsthead{
      \toprule
    }
    \tablehead{
      \multicolumn{#3}{l}{\small\sl continúa desde la página anterior}\\
      \hline
    }
    \tabletail{
      \hline
      \multicolumn{#3}{r}{\small\sl continúa en la página siguiente}\\
    }
    \tablelasttail{
      \hline
    }
    \bottomcaption{#1}
    \rowcolors[]{1}{cgoExtralight}{cgoLight}

  \begin{xtabular}{#2}
    #5
   \bottomrule
  \end{xtabular}
  \label{tabla:#4}
  \end{center}
}



\graphicspath{ {./img/} }

% Capítulos
\chapterstyle{bianchi}
\newcommand{\capitulo}[2]{
	\setcounter{chapter}{#1}
	\setcounter{section}{0}
	\setcounter{figure}{0}
	\setcounter{table}{0}
	\chapter*{\thechapter.\enskip #2}
	\addcontentsline{toc}{chapter}{\thechapter.\enskip #2}
	\markboth{#2}{#2}
}

% Apéndices
\renewcommand{\appendixname}{Apéndice}
\renewcommand*\cftappendixname{\appendixname}

\newcommand{\apendice}[1]{
	%\renewcommand{\thechapter}{A}
	\chapter{#1}
}

\renewcommand*\cftappendixname{\appendixname\ }

% Formato de portada
\makeatletter
\usepackage{xcolor}
\newcommand{\tutor}[1]{\def\@tutor{#1}}
\newcommand{\course}[1]{\def\@course{#1}}
\definecolor{cpardoBox}{HTML}{E6E6FF}
\def\maketitle{
  \null
  \thispagestyle{empty}
  % Cabecera ----------------
\noindent\includegraphics[width=\textwidth]{img/cabecera}\vspace{1cm}%
  \vfill
  % Título proyecto y escudo informática ----------------
  \colorbox{cpardoBox}{%
    \begin{minipage}{.8\textwidth}
      \vspace{.5cm}\Large
      \begin{center}
      \textbf{TFG del Grado en Ingeniería Informática}\vspace{.6cm}\\
      \textbf{\LARGE\@title{}}
      \end{center}
      \vspace{.2cm}
    \end{minipage}

  }%
  \hfill\begin{minipage}{.20\textwidth}
    \includegraphics[width=\textwidth]{escudoInfor}
  \end{minipage}
  \vfill
  % Datos de alumno, curso y tutores ------------------
  \begin{center}%
  {%
    \noindent\LARGE
    Presentado por \@author{}\\ 
    en Universidad de Burgos --- \@date{}\\
    Tutores: \@tutor{}\\
  }%
  \end{center}%
  \null
  \cleardoublepage
  }
\makeatother

\newcommand{\nombre}{Teodoro Ricardo García Sánchez} %%% cambio de comando
\newcommand{\titulo}{Aplicación de análisis de reseñas basado en modelos de lenguaje grandes mediante prompt engineering.}
\newcommand{\tutores}{Dª Virginia Ahedo García y D. Jose Ignacio Santos Martín}
% Datos de portada
\title{\titulo}
\author{\nombre}
\tutor{\tutores}
\date{\today}

\begin{document}

\maketitle


\newpage\null\thispagestyle{empty}\newpage


%%%%%%%%%%%%%%%%%%%%%%%%%%%%%%%%%%%%%%%%%%%%%%%%%%%%%%%%%%%%%%%%%%%%%%%%%%%%%%%%%%%%%%%%
\thispagestyle{empty}


\noindent\includegraphics[width=\textwidth]{img/cabecera}\vspace{1cm}

\noindent \tutores, profesores del departamento de Ingeniería de Organización.

\noindent Exponen:

\noindent Que el alumno D. \nombre, con DNI 50450469S, ha realizado el Trabajo final de Grado en Ingeniería Informática titulado \textbf{\titulo}. 

\noindent Y que dicho trabajo ha sido realizado por el alumno bajo la dirección de los que suscriben, en virtud de lo cual se autoriza su presentación y defensa.

\begin{center} %\large
En Burgos, {\large \today}
\end{center}

\vfill\vfill\vfill

% Author and supervisor
\begin{minipage}{0.45\textwidth}
\begin{flushleft} %\large
Vº. Bº. del Tutor:\\[2cm]
Dª Virginia Ahedo García
\end{flushleft}
\end{minipage}
\hfill
\begin{minipage}{0.45\textwidth}
\begin{flushleft} %\large
Vº. Bº. del co-tutor:\\[2cm]
D. Jose Ignacio Santos Martín
\end{flushleft}
\end{minipage}
\hfill

\vfill

% para casos con solo un tutor comentar lo anterior
% y descomentar lo siguiente
%Vº. Bº. del Tutor:\\[2cm]
%D. nombre tutor


\newpage\null\thispagestyle{empty}\newpage




\frontmatter

% Abstract en castellano
\renewcommand*\abstractname{Resumen}
\begin{abstract}
Este proyecto creará una aplicación para el análisis de sentimiento en las reseñas de usuario de productos y servicios.
\end{abstract}

\renewcommand*\abstractname{Descriptores}
\begin{abstract}
Análisis de sentimiento, prompt engineering, flask, python, vue, chatgpt, api rest.
\end{abstract}

\clearpage

% Abstract en inglés
\renewcommand*\abstractname{Abstract}
\begin{abstract}
This project will create an aplication for sentiment analysis in user reviews for diferent products and services.
\end{abstract}

\renewcommand*\abstractname{Keywords}
\begin{abstract}
Sentiment analysis, prompt engineering, flask, python, vue, chatgpt, api rest.
\end{abstract}

\clearpage

% Indices
\tableofcontents

\clearpage

\listoffigures

\clearpage

\listoftables
\clearpage

\mainmatter
\capitulo{1}{Introducción}

El análisis de sentimientos intenta descubrir la actitud de un usuario 
con respecto a algún tema.
Esta puede ser una reseña, un estado afectivo o emocional. 
Este de análisis se denomina minería de opinión.
El procesamiento de lenguaje natural busca mecanismos 
para la comunicación entre personas y máquinas por medio del
lenguaje natural. 
Se pueden realizar distintos tipos de análisis: 
sintáctico, semántico, prágmatico y planificación de frases.
En principio el objetivo del análisis de sentimientos es 
clasificar la polaridad del texto. (positiva, negativa o neutra).
Existen muchas formas de llevar a cabo este análisis (localización de palabras, 
afinidad léxica, métodos estadísticos o tecnicas conceptuales)
En este trabajo vamos a evaluar la utilización del "Prompt Engineering" 
asociado a un modelo grande de lenguaje (LLM) para ese análisis de sentimiento.
Evaluaremos su viabilidad y sus resultados. 
Usaremos dos medidas para evaluar el modelo, 
por un lado la precisión (número de elementos entre el número total) 
y la exhaustividad (recall) que mide la cantidad de elementos 
que ha sido capaz de identificar 
(se calcula dividiendo el número de aciertos entre el número 
total que hemos clasificado de esa forma). 
Se considera un modelo equiparable a un humano si consigue un 70 por ciento de aciertos.

\section{Estructura de la memoria}
\begin{description}
	\item[Introducción:] Descripción del tema sobre el que trata el trabajo y su situación. Incluye su estructura (memoria y anexos)
	\item[Objetivos del trabajo:] Este apartado explica de forma precisa y concisa cuales son los objetivos que se persiguen con la realización del proyecto. Se puede distinguir entre los objetivos marcados por los requisitos del software a construir y los objetivos de carácter técnico que plantea a la hora de llevar a la práctica el proyecto.
    \item[Conceptos teóricos:] Exposición de conceptos que facilitan la comprensión del proyecto.
    \item[Técnicas y herramientas:] Listado de metodologías y herramientas que han sido utilizadas para llevar a cabo el proyecto.
    \item[Aspectos relevantes:] Muestra aspectos a destacar durante la realización del proyecto.
    \item[Trabajos relacionados:] Estado del arte en el ámbito del sentiment analysis y trabajos similares.
    \item[Conclusiones y líneas de trabajo futuras:] Conclusiones obtenidas al final del proyecto y posibles ideas futuras.
\end{description}
\section{Estructura de los anexos}
\begin{description}
	\item[Plan de proyecto software:] Planificación temporal y viabilidad económica y legal.
	\item[Especificación de requisitos:] Objetivos y y requisitos establecidos al comienzo del proyecto.    
    \item[Especificación de diseño:] Recoge los diseños de datos, procedimental, arquitectónico y de interfaces.
    \item[Manual del programador:] Explica los conceptos más técnicos del proyecto como su instalación, la organización de carpetas y la ejecución.
    \item[Manual de usuario:] Es la guía de cómo utilizar la aplicación paso a paso.
    \item[Sostenibilización curricular:] Reflexión personal sobre los aspectos de la sostenibilidad que se abordan en el trabajo.
    
\end{description}

\capitulo{2}{Objetivos del proyecto}

El objetivo del trabajo propuesto es desarrollar una aplicación 
basada en inteligencia artificial (IA) utilizando modelos de lenguaje grandes (LLM~\cite{mirchandani2023large}) 
para el análisis, clasificación y calificación de reseñas de usuarios en portales 
de recomendación y tiendas de comercio electrónico. 
Los objetivos específicos son los siguientes:
\begin{itemize}
    \item[-] Crear una aplicación que utilice modelos de lenguaje grandes (como \href{https://chat.openai.com}{ChatGPT}, \href{https://bard.google.com/?hl=es}{Bard} o \href{https://research.facebook.com/publications/llama-open-and-efficient-foundation-language-models}{Llama})
    de manera efectiva para analizar y comprender las reseñas de usuarios. 
    \item[-] Diseñar un sistema basado en "prompts", definiendo las tareas, entradas y 
    salidas de la respuesta a una instrucción, para lograr una solución 
    ágil sin necesidad de nuevos entrenamientos.
    \item[-] Identificar y categorizar aspectos específicos del producto o servicio:
    \begin{description}
        \item[-] Implementar un sistema que pueda identificar y categorizar 
    los aspectos específicos mencionados por los usuarios en sus reseñas.
        \item[-] Extraer información relevante sobre productos o servicios, 
    identificando características específicas que generan opiniones.
    \end{description}
    \item[-] Relacionar opiniones y emociones con calificaciones:
    \begin{description}
        \item[-] Desarrollar un mecanismo que pueda relacionar las opiniones expresadas 
        en las reseñas con las calificaciones otorgadas por los usuarios.
        \item[-] Analizar las emociones asociadas a las opiniones para comprender 
        mejor la percepción de los usuarios respecto a productos y servicios.
    \end{description}
    \item[-] Implementar técnicas de ``\emph{prompt engineering}'' ~\cite{white2023prompt}:
    \begin{description}
        \item[-] Utilizar técnicas de ``\emph{prompt engineering}'' para optimizar la 
        interacción con el modelo de lenguaje, mejorando la calidad y 
        relevancia de las respuestas generadas.
        \item[-] Facilitar un desarrollo e implementación rápida y de bajo código para la solución propuesta. 
    \end{description}
        
\end{itemize}

El objetivo general es aprovechar la capacidad de los LLM 
para comprender el lenguaje natural y aplicarlos de manera efectiva en el 
análisis de reseñas de usuarios, proporcionando una solución ágil y 
eficiente para extraer información valiosa sobre la 
percepción de los usuarios en portales de recomendación y tiendas en línea.
\capitulo{3}{Conceptos teóricos}

Para una mejor compensión del trabajo realizado se recomienda conocer los conceptos teóricos mencionados en esta sección.

\section{Análisis de sentimiento}
El análisis de sentimiento es un campo de estudio en el procesamiento del lenguaje natural (PLN) que se centra en determinar y extraer la polaridad emocional asociada con un conjunto de datos de texto~\cite{Bedrava2022}. 

El objetivo principal es identificar y evaluar la actitud emocional expresada en un fragmento de texto, ya sea positiva, negativa o neutral.

El análisis de sentimiento implica varios pasos:
\begin{itemize}
	\item \textbf{Preprocesamiento del Texto}: antes de realizar el análisis de sentimiento, es necesario realizar una serie de tareas de preprocesamiento del texto, que incluyen la eliminación de palabras irrelevantes (stop words), lematización (reducción de palabras a su forma base), y la normalización del texto.
	\item \textbf{Representación del Texto}: el texto debe representarse de una manera que pueda ser procesada por algoritmos de aprendizaje automático. Esto a menudo implica convertir el texto en vectores numéricos utilizando técnicas como la bolsa de palabras (bag-of-words) o modelos de incrustación de palabras (word embeddings).
	\item \textbf{Modelado}: se utilizan modelos de aprendizaje automático, como máquinas de soporte vectorial (SVM), redes neuronales, o algoritmos de clasificación, para entrenar un modelo sobre los datos de texto anotados con etiquetas de sentimiento. Durante el entrenamiento, el modelo aprende patrones y características asociadas con la polaridad del sentimiento.
	\item \textbf{Clasificación}: una vez entrenado, el modelo se utiliza para predecir la polaridad del sentimiento de nuevos textos. El resultado típicamente indica si el texto es positivo, negativo o neutral.
\end{itemize}

El análisis de sentimiento se aplica en una variedad de campos, 
desde la minería de opiniones en redes sociales hasta la 
evaluación de comentarios de clientes en línea. 

Puede ser utilizado para comprender la percepción pública 
de productos, servicios o eventos, así como para tomar 
decisiones informadas en marketing, atención al cliente y 
otras áreas empresariales.

\section{Procesamiento del Lenguaje natural}
Es una rama de la inteligencia artificial que se ocupa de la interacción 
entre las computadoras y el lenguaje humano. \\
El objetivo principal del PLN es permitir que las máquinas comprendan, 
interpreten y generen texto de manera similar a como lo haría un ser humano~\cite{Moreno2018}.\\
Dentro del procesamiento del lenguaje natural hay una serie de conceptos que es importante conocer:
\begin{itemize}
	\item \textbf{Tokenización}: es el proceso de dividir un texto en unidades más 
	pequeñas llamadas "tokens". Un token puede ser una palabra, 
	una frase o incluso un carácter, dependiendo de la granularidad deseada.
	
	\item \textbf{Análisis morfológico}: implica comprender la estructura y la forma 
	de las palabras en un texto.\\ Esto puede incluir la descomposición 
	de las palabras en raíces y afijos para comprender su significado.

	\item \textbf{Análisis sintáctico}: se refiere a analizar la estructura 
	gramatical de las oraciones para comprender la relación entre las palabras.\\ 
	Esto implica identificar partes del discurso, como sustantivos, verbos, 
	adjetivos, etc., y cómo se combinan para formar oraciones 
	gramaticalmente correctas.
	
	\item \textbf{Análisis semántico}: busca comprender el significado de las palabras y 
	cómo se combinan para formar significados más complejos. \\
	Involucra la interpretación del contexto y la comprensión 
	de las relaciones entre las palabras.
	
	\item \textbf{Análisis pragmático}: considera el uso del lenguaje en situaciones 
	específicas y comprende el significado más allá de las palabras. \\
	Incluye la interpretación de intenciones, tono y contexto cultural.
	
	\item \textbf{Reconocimiento de entidades}: consiste en identificar y clasificar 
	entidades como nombres de personas, organizaciones, fechas, ubicaciones, etc., 
	dentro de un texto.
	
	\item \textbf{Desambiguación}: aborda la resolución de ambigüedades en el lenguaje.\\
	Por ejemplo, entender a qué se refiere un pronombre o resolver el significado 
	de una palabra con múltiples interpretaciones.
	
	\item \textbf{Generación de lenguaje}: implica la creación de texto humano legible por parte 
	de las máquinas.\\ 
	Puede ser tan simple como completar una oración o tan complejo 
	como la creación de contenido original.
		 
\end{itemize}


El procesamiento del lenguaje natural se utiliza en una variedad de 
aplicaciones, como motores de búsqueda, traducción automática, chatbots, 
resumen automático de textos, análisis de sentimientos y 
muchas otras áreas donde se requiere comprensión y generación de texto.\\


\section{Prompt Engineering}
El ``\emph{Prompt Engineering}'' se refiere a la práctica de diseñar de manera
 estratégica las instrucciones que se proporcionan a modelos de lenguaje, 
 como los ``\emph{Large Language Models}'', para obtener resultados específicos. 

En otras palabras, implica la formulación precisa de las solicitudes al modelo 
con el objetivo de influir en la salida de manera deseada~\cite{Solano2023}.

Esta técnica es especialmente relevante en el contexto de modelos 
generativos de lenguaje, como GPT-3, donde la calidad y la 
relevancia de las respuestas pueden depender en gran medida 
de cómo se presenta la información al modelo.\\
El ``\emph{Prompt Engineering}'' se utiliza para optimizar la interacción 
con estos modelos y obtener resultados que sean más útiles 
para una tarea particular.

Algunas estrategias comunes de ``\emph{Prompt Engineering}'' incluyen~\cite{white2023prompt}:
\begin{itemize}
	\item \textbf{Ajuste del Contexto Inicial}: proporcionar un contexto inicial específico que oriente al modelo hacia la tarea deseada.
	\item \textbf{Modificación de la Formulación}: experimentar con la redacción de las instrucciones para obtener respuestas más precisas o informativas.
	\item \textbf{Control de la Temperatura}: en modelos de lenguaje como GPT-3, se puede ajustar la "temperatura" para influir en la aleatoriedad de las respuestas generadas.
	\item \textbf{Incorporación de Ejemplos}: proporcionar ejemplos específicos dentro de la solicitud para guiar al modelo hacia comportamientos deseados.
\end{itemize}
El ``\emph{Prompt Engineering}'' es una parte importante de trabajar con modelos de lenguaje generativos para adaptar su salida a las necesidades específicas de una tarea o aplicación. 
La investigación continua en esta área busca mejorar la eficacia y la coherencia de la interacción con estos modelos.

\section{Patrones de diseño}
Los patrones de diseño en el ámbito del desarrollo de software 
son soluciones probadas para problemas recurrentes que surgen 
durante el diseño y la implementación de software.\\ 
Estos patrones representan las mejores prácticas y experiencias 
recopiladas por desarrolladores a lo largo del tiempo y 
proporcionan soluciones estructuradas y reutilizables 
para situaciones comunes en el desarrollo de software~\cite{fowler2012patterns}.\\
Los patrones de diseño pueden clasificarse en varias categorías, 
y algunos de los más conocidos son:
\begin{itemize}
	\item \textbf{Patrones de Creación}:
	\begin{itemize}
		\item \textbf{Singleton}: garantiza que una clase tenga solo una instancia y proporciona un punto global de acceso a ella.
		\item \textbf{Factory Method}: define una interfaz para crear un objeto, pero deja que las subclases alteren el tipo de objetos que se crearán.
	\end{itemize}
	\item \textbf{Patrones de Estructura}:
	\begin{itemize}
		\item \textbf{Adapter}: permite que la interfaz de una clase sea utilizada como otra interfaz.
		\item \textbf{Decorator}: añade funcionalidades a un objeto dinámicamente.
	\end{itemize}
	\item \textbf{Patrones de Comportamiento}:
	\begin{itemize}
		\item \textbf{Observer}: define una dependencia uno a muchos entre objetos para que, cuando un objeto cambie de estado, todos sus dependientes sean notificados y actualizados automáticamente.
		\item \textbf{Strategy}: define una familia de algoritmos, encapsula cada uno de ellos y los hace intercambiables.
	\end{itemize}
	\item \textbf{Patrones de Arquitectura}:
	\begin{itemize}
		\item \textbf{MVC (Modelo-Vista-Controlador)}: divide una aplicación en tres componentes interconectados, separando la representación de la información, la lógica de negocio y la gestión de la interfaz de usuario.
		\item \textbf{MVP (Modelo-Vista-Presentador)}: similar al MVC, pero con una separación más clara entre la lógica de presentación y la lógica de negocio.
	\end{itemize}
\end{itemize}

La aplicación de patrones de diseño puede llevar a un código más flexible, mantenible y escalable. Sin embargo, es importante aplicarlos con prudencia y adaptarlos según las necesidades específicas del proyecto, ya que una implementación incorrecta puede llevar a complicaciones innecesarias. Los patrones de diseño son herramientas poderosas, pero no son una solución única para todos los problemas.


\capitulo{4}{Técnicas y herramientas}

En esta parte de la memoria voy a presentar algunas de las técnicas de desarrollo 
que he utilizado para el desarrollo del proyecto, así como la información relevante 
de cada una de ellas.

\section{Técnicas}

\subsection{ORM~\cite{fowler2012patterns}}
\label{ORM}
Un ORM (Object-Relational Mapping) es una técnica de programación que se 
utiliza en el desarrollo de software para conectar objetos en un programa 
con tablas en una base de datos relacional. \\
Su propósito principal es permitir a los desarrolladores trabajar con datos de una base de datos relacional utilizando objetos y clases en lugar de escribir consultas SQL directamente. Un ORM mapea las estructuras de datos en una base de datos a objetos en un lenguaje de programación, lo que facilita la interacción con la base de datos y el manejo de datos de una manera más orientada a objetos.
Las ventajas clave de usar un ``\emph{ORM}'' incluyen:
\begin{itemize}
    \item \textbf{Abstracción de la base de datos}: Oºculta los detalles de la base de datos subyacente, 
    permitiendo a los desarrolladores interactuar con los datos en términos de objetos y clases en lugar de preocuparse por la sintaxis SQL específica de la base de datos.
    \item \textbf{Productividad}: Reduce la cantidad de código SQL que los desarrolladores deben escribir, 
    lo que ahorra tiempo y esfuerzo en el desarrollo de aplicaciones.
    \item \textbf{Portabilidad}: Puede ser diseñado para ser independiente de la base de datos, 
    lo que facilita la migración de una base de datos a otra sin cambiar el código de la aplicación.
    \item \textbf{Mantenibilidad}: Al trabajar con objetos y clases en lugar de consultas SQL directas, 
    el código puede ser más fácil de mantener y entender.
    \item \textbf{Seguridad}: Suelen ofrecer características de seguridad integradas, 
    como la prevención de inyecciones SQL, lo que reduce la vulnerabilidad de las aplicaciones a ataques. 
\end{itemize}
Ejemplos de ORMs populares incluyen Hibernate para Java, Entity Framework para .NET, 
SQLAlchemy para Python, Sequelize para Node.js y Django ORM para aplicaciones web Django en Python. \\
Cada uno de estos ORMs proporciona una forma de mapear objetos a tablas de base de datos y facilita 
la manipulación de datos en aplicaciones de software de manera más eficiente y mantenible.

\subsection{MVC~\cite{fowler2012patterns}}
MVC es un patrón arquitectónico comúnmente utilizado en el desarrollo de aplicaciones de software, 
especialmente en aplicaciones web. \\
Las siglas ``\emph{MVC}'' representa ``\emph{Model-View-Controller}'', y se divide en tres componentes principales 
que desempeñan roles específicos en la organización y estructura de una aplicación. 
Estos son:
\begin{itemize}
    \item \textbf{Model (Modelo)}: El modelo representa la lógica empresarial y los datos subyacentes 
    de la aplicación.\\ En otras palabras, el modelo se encarga de manejar la manipulación y 
    el acceso a los datos, así como de realizar cálculos y procesos necesarios.\\ 
    El modelo no tiene conocimiento de la interfaz de usuario y no se preocupa por cómo se muestran 
    los datos.\\ Es independiente de la vista y el controlador.
    \item \textbf{View (Vista)}: La vista se encarga de la presentación y la interfaz de usuario. \\
    Su función principal es mostrar los datos al usuario de una manera comprensible y atractiva. \\
    La vista no realiza operaciones comerciales o de lógica, sino que simplemente muestra la 
    información proporcionada por el modelo. \\
    En sistemas web, esto a menudo se traduce en la generación de páginas HTML o 
    interfaces de usuario dinámicas.
    item \textbf{Controller (Controlador)}: El controlador actúa como un intermediario entre el modelo y la vista. 
\end{itemize}

Recibe las solicitudes del usuario a través de la vista, procesa esas solicitudes, 
realiza las operaciones necesarias en el modelo y luego actualiza la vista para mostrar los resultados.\\
El controlador maneja la lógica de interacción y toma decisiones sobre qué acción realizar 
en función de la entrada del usuario.\\
La arquitectura MVC promueve una clara separación de preocupaciones, lo que facilita la modularidad y la escalabilidad de las aplicaciones. Esto permite a los desarrolladores trabajar de manera más eficiente, ya que pueden centrarse en una capa (modelo, vista o controlador) sin preocuparse demasiado por las otras capas. Además, el MVC es ampliamente utilizado en marcos de desarrollo web como Ruby on Rails, Django (Python), Laravel (PHP), y ASP.NET MVC (CSharp), entre otros.
En resumen, el patrón MVC es una forma efectiva de organizar y estructurar aplicaciones de software, lo que resulta en un código más claro, mantenible y flexible.

\subsection{LLM~\cite{mirchandani2023large}}
\label{LLM}
Un LLM (Modelo de Lenguaje Grande) es un tipo de modelo de inteligencia artificial
 diseñado para procesar y generar texto en lenguaje natural a gran escala.\\
Estos modelos están entrenados en enormes cantidades de datos de texto 
y están diseñados para comprender y generar texto en varios idiomas, 
lo que les permite llevar a cabo tareas de procesamiento de lenguaje natural de manera efectiva. \\
Algunos ejemplos son GPT-3, GPT-4 y BERT~\cite{BertSentimentModel1}.\\
Características clave:
\begin{itemize}
    \item \textbf{Tamaño y complejidad}: los modelos de lenguaje grandes contienen miles de millones o 
    incluso decenas de miles de millones de parámetros, lo que les permite capturar una 
    gran cantidad de conocimiento lingüístico y patrones de lenguaje.
    \item \textbf{Pre-entrenamiento y afinación}: estos modelos se entrenan primero en grandes corpus de texto, lo que les permite aprender el lenguaje natural y su estructura. Luego, se pueden afinar o ajustar para tareas específicas, como traducción automática, resumen de texto, generación de texto creativo, entre otros.
    \item \textbf{Generación de texto}: son capaces de generar texto coherente y contextualmente relevante.\\
    Pueden responder preguntas, completar oraciones, generar contenido original 
    y realizar tareas de procesamiento de lenguaje natural.
    \item \textbf{Amplia aplicabilidad}: estos modelos se utilizan en una amplia gama de aplicaciones, 
    que van desde asistentes virtuales y chatbots hasta motores de búsqueda mejorados, 
    análisis de sentimientos, resumen automático de texto, traducción automática y más.
\end{itemize}
 
Plantean desafíos éticos y de seguridad, ya que su 
capacidad para generar texto implica la necesidad de prevenir el uso indebido, 
la desinformación y el sesgo en el contenido generado.\\
Estos modelos han sido desarrollados y perfeccionados por empresas de tecnología y 
organizaciones de investigación en inteligencia artificial, y han demostrado 
un gran potencial en una variedad de aplicaciones. \\
A medida que continúan evolucionando, su influencia 
en la forma en que interactuamos con la tecnología y generamos contenido 
en línea también está en constante crecimiento.


\section{Herramientas}

\subsection{Python~\cite{lutz2001programming}}
Es un lenguaje de programación de alto nivel, interpretado y generalmente 
considerado como fácil de aprender y de utilizar. \\
Python destaca por su sintaxis clara y fácil de entender, 
lo que lo convierte en un lenguaje ideal tanto para principiantes 
como para programadores experimentados.\\
Es compatible con una amplia gama de sistemas operativos, 
lo que lo hace muy portable.\\

Ofrece una extensa biblioteca estándar que cubre una variedad de tareas, 
desde manipulación de archivos y redes hasta desarrollo web y 
matemáticas, lo que acelera el desarrollo de aplicaciones.\\
Cuenta con una comunidad de desarrollo muy activa y 
una gran cantidad de bibliotecas y marcos de trabajo de código abierto disponibles, 
lo que facilita la creación de aplicaciones complejas.\\
Es orientado a objetos, lo que significa que se enfoca en la creación de objetos 
que pueden contener datos y funciones.

Se utiliza en una amplia variedad de campos, como desarrollo web, 
análisis de datos, inteligencia artificial, automatización de tareas, desarrollo de juegos y más.

\subsection{Flask~\cite{dwyer2017flask}}
Flask es un framework de desarrollo web en Python que se utiliza 
para crear aplicaciones web de manera sencilla y rápida. \\
Está dentro de lo que se denomina ``\emph{microframework}'' por su enfoque minimalista, proporciona 
las funcionalidades esenciales pero permite agregar las bibliotecas necesarias según los requisitos.\\
No impone estructuras rígidas ni componentes innecesarios pero incorpora funcionalidades 
(como el motor de plantillas Jinja2 o el sistema de rutas) que hacen 
más sencillo el desarrollo de aplicaciones. \\
En este proyecto lo utilizamos para realizar un \href{https://aws.amazon.com/es/what-is/restful-api}{API RESTful}, que es un uso bastante común de este framework.


\subsection{SqlAlchemy~\cite{bayer2012sqlalchemy}}
SQLAlchemy es una biblioteca de Python que proporciona un conjunto de herramientas y
 un ORM(\ref{ORM}) que permite a los desarrolladores interactuar 
con bases de datos relacionales de una manera más sencilla y orientada a objetos.\\

Ofrece una amplia gama de métodos y operadores para realizar consultas en la base de datos, 
lo que facilita la construcción de consultas complejas de manera programática.

Gestiona las transacciones de manera transparente, lo que asegura la integridad 
de los datos al realizar operaciones de lectura y escritura.\\
Se pueden definir relaciones entre objetos y tablas de base de datos de manera sencilla,
 lo que simplifica la representación de datos relacionados.\\
Proporciona herramientas para realizar migraciones de bases de datos, 
lo que facilita la actualización de la estructura de la base de datos a medida 
que evoluciona la aplicación.\\
Permite la creación de extensiones y complementos personalizados 
para adaptarse a necesidades específicas.\\
Se utiliza principalmente en aplicaciones web, 
especialmente en el desarrollo de aplicaciones basadas en frameworks como Flask y Django. \\
Permite a los desarrolladores trabajar con bases de datos de manera más eficiente y 
mantener un código más limpio y legible al proporcionar una capa de abstracción entre la 
aplicación y la base de datos subyacente.\\ 
Esto también facilita la portabilidad de la aplicación entre diferentes sistemas de 
gestión de bases de datos.\\


\subsection{LlamaCpp~\cite{metallamaLM}}
La división de inteligencia artificial de Meta, llamada Meta AI, presentó LLaMA, un LLM (\ref{LLM}) de 
65.000 millones de parámetros que permite disfrutar de un motor de 
IA conversacional muy parecido a ChatGPT.\\
Este modelo estaba inicialmente disponible para desarrolladores e investigadores 
que justificaran su uso.\\
Grigori Gerganov publicó en Github~\cite{gerganovllamaCpp} un pequeño desarrollo llamado llama.cpp,
un proyecto que permite poder usar el modelo LLaMA en un portátil o un PC convencional.\\ 
Eso se logra gracias a la llamada ``\emph{cuantización}'' que reduce el tamaño de los modelos de 
Facebook para hacerlos ``manejables'' por equipos más modestos a nivel de hardware.

\subsection{PostgreSql~\cite{juba2015learning}}
PostgreSQL es un sistema de gestión de bases de datos relacional de código abierto. 
Admite consultas complejas y procedimientos almacenados.\\
Proporciona mecanismos para garantizar la integridad de los datos almacenados, 
como restricciones de clave primaria y foránea, así como comprobaciones de restricciones.\\
Es escalable y se puede utilizar en aplicaciones de todos los tamaños, 
desde aplicaciones pequeñas hasta sistemas empresariales de alto rendimiento.\\
Admite la replicación y puede configurarse para lograr alta disponibilidad mediante 
la implementación de réplicas y clústeres.
Soporta autenticación, autorización y cifrado de datos.\\
PostgreSQL es una opción popular tanto en la comunidad de código abierto como 
en empresas que buscan una base de datos de alto rendimiento y confiabilidad para sus aplicaciones.\\
Su licencia de código abierto permite su uso y distribución sin costos de licencia, 
lo que lo convierte en una opción atractiva para una amplia variedad de proyectos.

\subsection{Docker~\cite{garzas2105docker}}
Docker es una plataforma que permite empaquetar una aplicación y 
todas sus dependencias en un contenedor.\\
Este contenedor es como una unidad autónoma que puede ejecutarse de manera 
consistente en cualquier entorno que admita Docker. 
Además, Docker facilita la gestión y la implementación de estas aplicaciones empaquetadas.
Al utilizar Docker, los desarrolladores pueden estar seguros de que la aplicación 
se ejecutará de la misma manera en todas partes, desde el entorno de desarrollo hasta el de producción. 
Esto elimina problemas relacionados con las diferencias entre configuraciones y facilita las pruebas.
Se pueden iniciar y detener rápidamente, lo que permite una implementación más rápida de 
nuevas versiones o características. \\
Esto es crucial en el \emph{desarrollo ágil}, donde la velocidad de respuesta a los cambios es esencial.

Además, encapsula todas las dependencias de la aplicación en el contenedor, 
evitando conflictos con otras aplicaciones o componentes del sistema. 
Esto simplifica la gestión de dependencias y reduce los posibles problemas de compatibilidad. \\
También facilita la escalabilidad horizontal, lo que significa que puedes 
ejecutar múltiples instancias de contenedores de la misma aplicación para 
gestionar cargas de trabajo más pesadas de manera eficiente.

Los contenedores Docker pueden compartirse fácilmente, lo que facilita la colaboración 
entre equipos de desarrollo y operaciones. 
Todos trabajan con la misma configuración, lo que reduce la posibilidad de errores 
causados por diferencias en los entornos.

\subsection{ChatGPT~\cite{chatgpt1}}
ChatGPT es un modelo de lenguaje desarrollado por OpenAI, 
basado en la arquitectura GPT (Generative Pre-trained Transformer). 

La arquitectura GPT es una red neuronal de tipo transformer que ha sido 
preentrenada en grandes cantidades de datos textuales para aprender patrones y 
estructuras del lenguaje.

ChatGPT puede generar respuestas contextualmente relevantes en función de las 
entradas que recibe. Puede ser utilizado para una variedad de aplicaciones, 
desde responder preguntas y ayudar en la redacción de textos.

Es importante tener en cuenta que ChatGPT no posee conocimiento del mundo 
en tiempo real ni tiene la capacidad de razonamiento 
ni comprensión como lo haría un ser humano.

Durante la realización del trabajo fin de grado he usado chatgpt para obtener 
información y sugerencias siempre contrastándolas con otras fuentes para comprobar 
su corrección.

Es especialmente interesante su uso para hacer resúmenes y organizar la información 
de forma coherente.

\capitulo{5}{Aspectos relevantes del desarrollo del proyecto}

Este apartado pretende recoger los aspectos más interesantes del desarrollo del proyecto, comentados por los autores del mismo.
Debe incluir desde la exposición del ciclo de vida utilizado, 
hasta los detalles de mayor relevancia de las fases de análisis, diseño e implementación.
Se busca que no sea una mera operación de copiar y pegar diagramas y extractos del código fuente, 
sino que realmente se justifiquen los caminos de solución que se han tomado, 
especialmente aquellos que no sean triviales.
Puede ser el lugar más adecuado para documentar los aspectos más interesantes del diseño y 
de la implementación, con un mayor hincapié en aspectos tales como el tipo de arquitectura elegido, 
los índices de las tablas de la base de datos, normalización y desnormalización, 
distribución en ficheros3, reglas de negocio dentro de las bases de datos (EDVHV GH GDWRV DFWLYDV), 
aspectos de desarrollo relacionados con el WWW...
Este apartado, debe convertirse en el resumen de la experiencia práctica del proyecto, y 
por sí mismo justifica que la memoria se convierta en un documento útil, 
fuente de referencia para los autores, los tutores y futuros alumnos.

\section{Metodologías}

\subsection{Metodologías ágiles(\href{https://www.scrum.org/}{SCRUM})}
Scrum es un marco de trabajo ágil utilizado comúnmente en el desarrollo de software, 
aunque se ha extendido a otras áreas. \\
Proporciona un enfoque estructurado para la gestión de proyectos que se centra en la entrega 
iterativa y incremental de productos. \\
Scrum se basa en los principios del manifiesto ágil y busca mejorar la eficiencia, 
la flexibilidad y la colaboración en equipos de desarrollo.

\subsubsection{Roles}

\begin{itemize}
   \item \textbf{Scrum Master}: Facilita el proceso Scrum, elimina obstáculos y ayuda al equipo a alcanzar sus objetivos.
   \item \textbf{Product Owner}: Representa las necesidades del cliente y define las características del producto.
   \item \textbf{Equipo de Desarrollo}: Profesionales que trabajan en la entrega del producto.
\end{itemize}

\subsubsection{Eventos}

\begin{itemize}
   \item \textbf{Sprint}: Un periodo de tiempo fijo (generalmente de 2 a 4 semanas) en el que se entrega un incremento de producto.
   \item \textbf{Reunión de Planificación del Sprint}: Al inicio de cada sprint, el equipo planifica el trabajo que se realizará durante ese periodo.
   \item \textbf{Revisión del Sprint}: Al final de cada sprint, el equipo presenta el trabajo completado al Product Owner y a otras partes interesadas.
   \item \textbf{Retrospectiva del Sprint}: Una sesión al final de cada sprint donde el equipo revisa su desempeño y busca formas de mejorar.
\end{itemize}

\subsubsection{Artefactos}

\begin{itemize}
\item \textbf{Product Backlog}: Una lista priorizada de todas las funcionalidades, cambios y mejoras propuestas para el producto.
\item \textbf{Sprint Backlog}: La lista de tareas que el equipo se compromete a completar durante un sprint.
\item \textbf{Incremento}: El producto funcional y potencialmente entregable al final de cada sprint.
\end{itemize}

Scrum promueve la transparencia, la inspección y la adaptación, lo que significa que los equipos pueden ajustar su enfoque y estrategia en función de los cambios en los requisitos del cliente o en las circunstancias del proyecto. Este marco de trabajo es especialmente útil en entornos donde los requisitos pueden cambiar con frecuencia y se valora la capacidad de respuesta y flexibilidad del equipo de desarrollo.

Para este proyecto se ha seguido la metodología \textbf{SCRUM} configurando los sprints con una duración de 2 semanas,
para la planificación se ha utilizado la herramienta \href{https://zube.io/tgs1003/tfg_2023_2024}{Zube}.
Se puede encontrar más detalles sobre el proceso seguido en los anexos de este documento.

\subsection{DevOps}
\textbf{DevOps} es una cultura, filosofía y conjunto de prácticas que se centran en la colaboración 
estrecha y la integración continua entre los equipos de desarrollo (Dev) y operaciones (Ops) 
en el ciclo de vida del desarrollo de software. \\
El objetivo principal de \textbf{DevOps} es acelerar la entrega de software, 
mejorar la calidad y la confiabilidad, y permitir una respuesta rápida a los 
cambios y a las necesidades de los usuarios. \\DevOps promueve la automatización, 
la comunicación eficaz y la colaboración entre los equipos, lo que permite un flujo 
de trabajo más eficiente en todo el proceso de desarrollo y entrega de software.\\
Algunos de los principios y prácticas clave de \textbf{DevOps} incluyen:
\begin{itemize}
   \item \textbf{Automatización}: La automatización de tareas repetitivas y procesos manuales acelera la entrega y minimiza los errores. Esto incluye la automatización de pruebas, implementaciones, aprovisionamiento de infraestructura y monitoreo.
   \item \textbf{Integración Continua (CI)}: Los cambios de código se integran regularmente en un repositorio compartido, 
se prueban automáticamente y se implementan en un entorno de desarrollo o de prueba. 
Esto asegura que el código esté siempre en un estado funcional.
\item \textbf{Entrega Continua (CD)}: La entrega continua extiende la integración continua al permitir 
la entrega automática de cambios a entornos de prueba o producción después de la integración y 
las pruebas exitosas.
\item \textbf{Monitoreo y Retroalimentación Continua}: El monitoreo constante de aplicaciones y sistemas 
permite detectar problemas en tiempo real y proporciona información valiosa para 
mejorar la calidad y la eficiencia.
\item \textbf{Colaboración y Comunicación}: La comunicación efectiva entre los equipos de desarrollo y 
operaciones es esencial. La colaboración se fomenta mediante reuniones regulares, 
herramientas compartidas y un entendimiento mutuo de las metas y responsabilidades.
\item \textbf{Infraestructura como Código (IaC)}: La infraestructura se administra y despliega como código, 
lo que facilita la creación y el mantenimiento de entornos de desarrollo y producción 
consistentes y escalables.
\item \textbf{Cultura de Mejora Continua}: DevOps promueve una cultura en la que se aprende 
de los errores y se busca constantemente la mejora en todos los aspectos del desarrollo y 
la operación de software.
\item \textbf{Seguridad}: La seguridad es un aspecto crítico de \textbf{DevOps}. 
Los principios de seguridad deben estar integrados en todas las etapas 
del ciclo de vida del desarrollo de software.
\end{itemize}
La implementación exitosa de \textbf{DevOps} puede llevar a una mayor velocidad de entrega, 
una mayor calidad del software, una mayor eficiencia operativa y 
una mayor capacidad de respuesta a los cambios del mercado y las necesidades del cliente. 
Esta metodología se ha vuelto esencial en el desarrollo de software moderno, 
especialmente en entornos ágiles y de entrega continua.

Para el publicar del código de este trabajo se ha utilizado \href{https://www.docker.com/}{Docker} junto con \href{https://www.portainer.io/}{Portainer} para el 
despliegue continuo.


\section{Desarrollo del proyecto}

El principal escollo que me he encontrado para la realización de este proyecto es la elección del modelo
LLM adecuado. \\
Por un lado tenemos el servicio de \href{https://openai.com}{OpenAI}. Este servicio es quizá el más estable y que 
mejores resultados obtiene.\\ 
Sin embargo este servicio es de pago (aunque propocionan un saldo gratuito al crear una cuenta, 
no permite realizar más de 3 llamadas al minuto).
Además tenemos el tiempo de proceso para cada reseña (unos 3 segundos).
Existe otros LLMs disponibles como por ejemplo \href{https://ai.meta.com/}{Llama2 (Meta)}, 
\href{https://bard.google.com/chat}{Bard(Google)}.\\
Con el modelo de Meta existe la posibilidad de ejecutarlo en un entorno local mediante 
diferentes librerías (LlamaCpp).\\
Otros modelos que se pueden ejecutar en local son \href{https://gpt4all.io/index.html}{GPT4All}.
Una herramienta muy útil para elegir el modelo adecuado es \href{https://lmstudio.ai/}{LM Studio}, permite descargar y probar distintos modelos,
además ofrece una API compatible con OpenAI.\\
Para la ejecución en local es muy recomendable el uso de una tarjeta gráfica 
para optimizar la ejecución del modelo.

Al planificar el segundo sprint incluí una tarea para crear un prototipo en collab, 
sin embargo esta tarea tenía unas dependencias que hacían imposible su realización.
Lo primero que tenía que hacer era seleccionar un dataset para poder realizar las pruebas.
Y además necesitaba seleccionar el modelo LLM idoneo para la realización de este protoripo.

\section{Elección del conjunto de datos}

Elegir un conjunto de datos adecuado es crucial para el análisis de sentimientos.
En este proyecto nos insteresa realizar un análisis de sentimientos en reseñas de 
Amazon por lo que se ha buscado con esa premisa. \\
Esto reduce la búsqueda pero aun así hay que encontrar 
un conjunto de datos que tenga una serie de características que lo haga ideal para su uso en el proyecto.\\
Por un lado, como veremos más adelante, el análisis de sentimiento tiene un coste, 
ya sea económico (OpenAI), o en tiempo (LlamaCpp).
Esto implica que tenemos que elegir un dataset de un tamaño adecuado.
Por un lado, tiene que ser suficientemente grande para que los resultados sean significativos y 
por otro lado lo suficientemente pequeño como para que el coste sea adecuado al proyecto.
El conjunto de datos debe contener información suficiente como para poder realizar 
comparaciones bajo distintas circustancias de estudio, por usuario, por producto, etc.\\
También sería interesante que el conjunto de datos sea represativo de la población, 
para ello lo ideal sería obtener el conjunto de datos de diferentes orígenes.\\
En nuestro caso el origen es siempre Amazon pero a la hora de seleccionar los datos 
intentaremos que estos sean lo más variados posibles.\\
Además hay tener en cuenta la disponibilidad, en algunos casos Amazon ha decidido anular 
la licencia de uso con lo que esos conjuntos de datos no se podrían usar.\\
En base a estas premisas se han estudiado datasets de diferentes orígenes:

\subsection{Universidad de California~\cite{jin018}}
La Universidad de California proporciona un conjunto de datos con reseñas de Amazon. Los más llamativo de este conjunto de datos es su tamaño (230 millones de reseñas) y la cantidad de metadatos incluidos.
También propocionan ejemplos de uso y de código así como un cuaderno en \href{https://colab.research.google.com/drive/1Zv6MARGQcrBbLHyjPVVMZVnRWsRnVMpV}{Colab}.
\par
\href{https://cseweb.ucsd.edu/~jmcauley/datasets/amazon_v2/}{Amazon Reviews Dataset}

\subsection{Kaggle}
Kaggle es una plataforma en línea que ofrece una variedad de recursos relacionados con la ciencia de datos y el aprendizaje automático. Fue fundada en 2010 y adquirida por Google en 2017. Kaggle proporciona un entorno donde los científicos de datos, los investigadores y los entusiastas del aprendizaje automático pueden colaborar, compartir conocimientos, participar en competiciones de ciencia de datos y acceder a conjuntos de datos.
Uno de los conjuntos de datos evaluados han sido:
\begin{itemize}
   \item \href{https://www.kaggle.com/datasets/bittlingmayer/amazonreviews}{Amazon Reviews for Sentiment Analysis}
   \par Este conjuntos de datos tiene un tamaño adecuado para nuestro proposito pero le faltan metadatos y no nos permitiría desarrollar toda la funcionalidad esperada en la aplicación.
   \item \href{https://www.kaggle.com/datasets/kritanjalijain/amazon-reviews}{Amazon Reviews}
\end{itemize}



\subsection{Hugging Face}
Hugging Face es una empresa y plataforma en línea que se centra en la creación y distribución de modelos de lenguaje de inteligencia artificial, así como en herramientas y recursos relacionados con el procesamiento del lenguaje natural (NLP). La compañía es conocida por su biblioteca de modelos preentrenados, así como por su contribución a la comunidad de aprendizaje automático y NLP.
Existen muchos conjuntos de datos que podemos usar para nuestro proyecto, por ejemplo:
\begin{itemize}
   \item \href{https://huggingface.co/datasets/mesmalif/amazon-shoe-reviews}{Amazon Shoe Reviews}
   \item \href{https://huggingface.co/datasets/LoganKells/amazon_product_reviews_video_games}{Amazon Video Games Review}
   \item \href{https://huggingface.co/datasets/hugginglearners/amazon-reviews-sentiment-analysis}{Amazon Reviews for Sentiment Analysis}   
\end{itemize}



\section{Elección del modelo LLM}

Para la elección del modelo a usar en este trabajo, primero examinamos el modelo más sencillo de utilizar.
Este es el modelo de OpenAI~\cite{chatgpt1}. En concreto el modelo \textbf{ChatGpt 3.5 Turbo}, este modelo es muy estable, rápido y económico.
El problema de usar este modelo, a pesar de ser económico, es su coste.
Para evitar ese problema y poder utilizar un modelo sin limitaciones, he buscado otras alternativas.
Entre estas se encuentran:
\begin{itemize}
   \item \textbf{Microsoft Azure Language Models} Microsoft Azure ofrece varios servicios de procesamiento del lenguaje natural, incluidos modelos de lenguaje como parte de su oferta de servicios cognitivos.
   \item \textbf{Google Cloud Natural Language API} Google Cloud proporciona una API de procesamiento del lenguaje natural que incluye funciones avanzadas como análisis de sentimientos, extracción de entidades y más.
   \item \textbf{Hugging Face Transformers} Hugging Face es una plataforma que ofrece acceso a una amplia variedad de modelos de lenguaje preentrenados. Los modelos GPT y otros están disponibles a través de su biblioteca Transformers.
   \item \textbf{Facebook AI} Facebook AI Research (FAIR) trabaja en el desarrollo de modelos de lenguaje y herramientas relacionadas con la inteligencia artificial. Aunque no tienen un modelo específico comparable a GPT, han contribuido significativamente al campo.
   \item \textbf{Rasa} Rasa es una plataforma de código abierto para construir asistentes conversacionales. Permite a los desarrolladores crear chatbots personalizados y asistentes virtuales.
   \item \textbf{Chatbot Frameworks} Hay varios marcos y bibliotecas de código abierto que permiten a los desarrolladores construir sus propios chatbots, como ChatterBot, Botpress, y Microsoft Bot Framework.
   \item \textbf{Dialogflow} Dialogflow, propiedad de Google, es una plataforma de desarrollo de chatbots y asistentes virtuales que utiliza tecnologías de procesamiento del lenguaje natural.
\end{itemize}

Investigando sobre las diferentes alternativas nos hemos encontrado con 
dificultades a la hora de elegir la idónea.
Por un lado, están los modelos que también son de pago como son ``Microsoft Azure Language Models'' 
y Google Cloud Natural Language API. Ambas soluciones ofrecen periodos de prueba limitados.
Por otro lado están las plataformas para asistentes conversacionales como Dialogflow, Rasa, etc 
que no ofrecen un API para poder usarlas en un programa.
Por último nos queda \textbf{Facebook API}, en este caso su modelo 
\href{https://github.com/facebookresearch/llama}{Llama2} está disponible para todo tipo de usos.
Existen una colección de modelos que van desde 7b a 70b (7b significa que se han usado 
7.000 millones de tokens para su entrenamiento y 70b 70.000 millones).\\
Aún así el modelo más pequeño necesita unos recursos muy elevados para su funcionamiento. 
Aquí es dónde entra en juego la \href{https://es.wikipedia.org/wiki/Cuantificaci%C3%B3n_digital}{cuantificación}.\\
Aplicando la cuantificación, los modelos ocupan menos espacio y son más rápidos pero también más imprecisos.
En la plataforma Hugging Face existen multitud de modelos de diferentes tamaños y tipos de cuantificación.
En este punto nos quedaba elegir un modelo que se pueda ejecutar de una manera razonable 
en equipos personales pero manteniendo unos niveles razonables de precisión.\\
Así entra en juego el último de los elementos para poder elegir un modelo y 
es un método sencillo para poder evaluar los modelos.\\
Para ello he usado una aplicación llamada \href{https://lmstudio.ai/}{LM Studio}.\\ 
Es una aplicación gratuita que permite descargar modelos de Hugging Face y 
probarlos mediante un chat o mediante una API.\\
Además permite aprovechar una tarjeta gráfica en el caso de tenerla 
para poder usar el modelo aprovechando sus capacidades de computación en paralelo.


\capitulo{6}{Trabajos relacionados}

\section{Lingmotif~\cite{lingmotif}}
Lingmotif es una aplicación multi-plataforma de sobremesa que analiza textos desde la perspectiva 
del Análisis de Sentimiento. \\Básicamente, es capaz de determinar la orientación semántica 
(si es positivo o negativo y en qué grado) de un texto o conjunto de textos, mediante \textbf{la detección 
de expresiones lingüísticas} que indican una determinada polaridad.

A diferencia de la mayoría del software existente, Lingmotif no es un sólo un clasificador, 
ya que no se limita a clasificar un texto como positivo o negativo, sino que además ofrece una 
serie de datos cuantitativos, una visualización del ``perfil de sentimiento'' del texto o textos 
(incluyendo series temporales), y un detallado análisis cualitativo del texto en sí, en el que 
se muestran los segmentos textuales identificados. \\Estas funcionalidades lo convierten en un herramienta 
única, y sus aplicaciones van más allá de las que normalmente ofrece este tipo de software. 
Lingmotif ofrece los resultados de sus análisis en archivos con formato HTML, con la versatilidad y 
fácil manejo que esto supone.

Actualmente Lingmotif analiza textos en español e inglés. Se está trabajando en nuevas versiones 
para alemán e italiano.

\section{\href{https://www.lexalytics.com/}{Lexalytics}}
Lexalytics es una empresa que se especializa en tecnologías de análisis de texto y 
procesamiento del lenguaje natural (PLN).\\
Ofrece soluciones para extraer información de datos de texto no estructurados, 
como contenido de redes sociales, comentarios de clientes, reseñas y más. 
Los productos de Lexalytics están diseñados para ayudar a las empresas a analizar 
y comprender el sentimiento, los temas y las tendencias dentro de grandes volúmenes de información textual.
Uno de sus productos destacados es la Plataforma de Inteligencia Lexalytics, 
que incluye varias herramientas y funciones para el análisis de texto. 
Tienen un API (al que llaman Semantria) que se usa para análisis de texto y sentimiento.
Admite varios idiomas, lo que las hace adecuadas para empresas con presencia global.\\
Estas herramientas y servicios pueden ser valiosos para empresas que 
buscan dar sentido a la gran cantidad de texto no estructurado que generan o 
encuentran en el curso de sus operaciones. \\Se pueden aplicar en áreas como la gestión 
de la experiencia del cliente, el monitoreo de redes sociales, la investigación de mercado 
y otros campos donde comprender la información textual es crucial. 

\section{\href{https://www.ibm.com/products/watson-studio}{IBM Watson Studio}}
IBM Watson Studio es una plataforma de inteligencia artificial (IA) desarrollada 
por IBM que proporciona servicios y herramientas para el procesamiento 
del lenguaje natural y otros. 
Lleva el nombre del fundador de IBM (Thomas J. Watson)
También utiliza técnicas de aprendizaje automático 
para extraer patrones a partir de grandes conjuntos de datos y realizar análisis predictivos.\\
Tiene capacidades para analizar y comprender imágenes y videos, lo que puede 
ser útil en aplicaciones como el reconocimiento facial y la interpretación de contenido visual.\\
Se ha utilizado en aplicaciones como juegos de preguntas y 
respuestas, demostrando su capacidad para competir y ganar contra humanos en contextos complejos.
Puede automatizar tareas y procesos empresariales, mejorando la eficiencia operativa.
Analiza datos en tiempo real para proporcionar información instantánea y tomar decisiones basadas en datos.
IBM Watson se ha aplicado en diversos campos obtenido muy buenos resultados.
Permite crear un cuanta gratuita (limitada) para poder probarla.

\section{Sentinel}
Sentinel~\cite{Sentinel1} es una aplicación web desarrollada como TFG del grado de ingeniería 
informática de la Universidad de Burgos que realiza un análisis de sentimientos de textos 
extraídos de las redes sociales Twitter e Instagram.\\
La aplicación busca los tweets relacionados con la palabra introducida 
en el caso de Twitter, o los comentarios que le han escrito a la 
cuenta de Instagram que se ha buscado.\\
Después analiza el sentimiento que hay en ellos, los puntúa con valores entre 0 y 1 y 
se almacenan los resultados.\\
Estos resultados se muestran al usuario en gráficos 
y tablas para hacer la experiencia más visual.\\ 
Además se le ofrece la opción de calcular series temporales a partir de los resultados, 
y se da una predicción de los valores futuros.

\section{\href{https://www.danielsoper.com/sentimentanalysis/default.aspx}{Free Sentiment Analysis}}
Es una herramienta gratuita que permite realizar un analisis de sentimiento de cualquier 
texto escrito en Inglés. La herramienta califica el sentimiento con un número entre -100 y 100 
dependiendo del sentimiento detectado en el texto.
Sólo hay que pegar el texto en el cuadro de texto y pulsar el botón ``Analyze text''

Para su funcionamiento utiliza algoritmos de linguística computacional y minería de texto.
El modelo se ha entrenado usando el ``American National Corpus'' con lo que sólo funciona con 
textos en inglés americano y escritos después de 1990. 
No se conoce su funcionamiento interno pero por la descripción parece que ha usado un 
corpus anotado y usa las coincidencias de palabras para dar una nota al texto dependiendo 
de las palabras encontradas.
\capitulo{7}{Conclusiones y Líneas de trabajo futuras}

Todo proyecto debe incluir las conclusiones que se derivan de su desarrollo. 
Éstas pueden ser de diferente índole, dependiendo de la tipología del proyecto, 
pero normalmente van a estar presentes un conjunto de conclusiones relacionadas 
con los resultados del proyecto y un conjunto de conclusiones técnicas. 
Además, resulta muy útil realizar un informe crítico indicando cómo se puede 
mejorar el proyecto, o cómo se puede continuar trabajando en la línea 
del proyecto realizado. 
\section{Conclusiones}


\section{Líneas de trabajo futuras}

\subsection{Nuevos modelos}

\subsection{Nuevas fuentes de datos}
Web service

\subsection{Más opciones de análisis}
Configuración del prompt

\subsection{GridFS}



\bibliographystyle{plain}
\bibliography{bibliografia}

\end{document}
